\documentclass[../libro]{subfiles}

\begin{document}

\ifSubfilesClassLoaded{\mainmatter\chapter{行列式}\clearpage}{}

\section{单位阵}

本节, 我要介绍一类特别的方阵, 单位阵.

\begin{definition}[\(\delta\)-记号]
    设 \(i\), \(j\) 是二个文字.
    定义
    \begin{align*}
        \delta(i, j)
        = \begin{cases}
              0, & i \neq j; \\
              1, & i = j.
          \end{cases}
    \end{align*}
\end{definition}

显然, \(\delta\)-记号是表示二个文字是否相等的一个量.
\(\delta(i, j) = 1\), 说明 \(i\), \(j\) 是同一个文字;
\(\delta(i, j) = 0\), 说明 \(i\), \(j\) 是不同的二个文字.

\begin{definition}[单位阵]
    设 \(n\) 为正整数.
    作 \(n\)~级阵 \(I_n\),
    使 \([I_n]_{i,j} = \delta(i, j)\),
    (\(1 \leq i, j \leq n\)).
    我们说, \(I_n\) 是 \(n\)~级\emph{单位阵}.

    当我们不强调单位阵的尺寸时,
    我们可简单地写单位阵为 \(I\).
\end{definition}

\(I\) 来自英语 \angla{identity}.

\begin{example}
    \(1\)~级单位阵是 \([1]\).
    \(2\)~级单位阵是
    \(
    \begin{bmatrix}
        1 & 0 \\
        0 & 1 \\
    \end{bmatrix}.
    \)
    \(3\)~级单位阵是
    \(
    \begin{bmatrix}
        1 & 0 & 0 \\
        0 & 1 & 0 \\
        0 & 0 & 1 \\
    \end{bmatrix}.
    \)

    一般地, \(n\)~级单位阵含 \(n\)~个 \(1\)
    与 \(n^2 - n\)~个 \(0\),
    且 \(1\) 恰在行号与列号相等的地方
    (也就是, \(0\) 恰在行号与列号不等的地方).
\end{example}

可写一个 \(n \times 1\)~阵%
为 \(n\)~级单位阵 \(I_n\)~的列的数乘的和.
具体地, 我们设 \(x\) 是一个 \(n \times 1\)~阵,
且设 \(I_n = [e_1, e_2, \dots, e_n]\)
(也就是,
设 \(e_1\), \(e_2\), \(\dots\), \(e_n\) 是
\(I_n\) 的列 \(1\), \(2\), \(\dots\), \(n\)).
记 \(k_i = [x]_{i,1}\).
则
\begin{align*}
    x
    % = {} &
    =
    k_1 e_1 + k_2 e_2 + \dots + k_n e_n
    % \\
    % = {} &
    =
    \sum_{\ell = 1}^{n} {k_\ell e_\ell}.
\end{align*}
为证明此式, 我们要用%
单位阵的定义、加法的定义、数乘的定义
(不难看出, 等式二侧的尺寸都是 \(n \times 1\)):
\begin{align*}
    \left[ \sum_{\ell = 1}^{n} {k_\ell e_\ell} \right]_{i,1}
    = {} &
    \sum_{\ell = 1}^{n} {[k_\ell e_\ell]_{i,1}}
    \\
    = {} &
    \sum_{\ell = 1}^{n} {k_\ell [e_\ell]_{i,1}}
    \\
    = {} &
    \sum_{\ell = 1}^{n} {k_\ell [I_n]_{i,\ell}}
    \\
    = {} &
    \sum_{\ell = 1}^{n} {k_\ell \delta(i, \ell)}
    \\
    = {} &
    k_i \delta(i, i)
    + \sum_{\substack{1 \leq \ell \leq n \\ \ell \neq i}}
    {k_\ell \delta(i, \ell)}
    \\
    = {} &
    k_i
    + \sum_{\substack{1 \leq \ell \leq n \\ \ell \neq i}}
    {k_\ell 0}
    \\
    = {} &
    k_i.
\end{align*}

至多以一个方式写一个 \(n \times 1\)~阵%
为 \(n\)~级单位阵 \(I_n\)~的列的数乘的和.
具体地, 设 \(e_1\), \(e_2\), \(\dots\), \(e_n\) 是
\(I_n\) 的列 \(1\), \(2\), \(\dots\), \(n\).
再设
\begin{align*}
    \sum_{\ell = 1}^{n} {k_\ell e_\ell}
    = \sum_{\ell = 1}^{n} {p_\ell e_\ell}.
\end{align*}
不难算出
\begin{align*}
     & \sum_{\ell = 1}^{n} {k_\ell e_\ell}
    = \begin{bmatrix}
          k_1    \\
          \vdots \\
          k_n
      \end{bmatrix},                       \\
     & \sum_{\ell = 1}^{n} {p_\ell e_\ell}
    = \begin{bmatrix}
          p_1    \\
          \vdots \\
          p_n
      \end{bmatrix}.
\end{align*}
从而
\begin{align*}
    \begin{bmatrix}
        k_1    \\
        \vdots \\
        k_n
    \end{bmatrix}
    = \begin{bmatrix}
          p_1    \\
          \vdots \\
          p_n
      \end{bmatrix}.
\end{align*}
由此可知, \(k_i = p_i\).

综上, 我们有

\begin{theorem}
    可以, 且只能以一个方式写一个
    \(n \times 1\)~阵为 \(n\)~级单位阵 \(I_n\)~的%
    列的数乘的和.
\end{theorem}

\begin{example}
    设
    \(
    x = \begin{bmatrix}
        6 \\ 2 \\ 8 \\
    \end{bmatrix}.
    \)
    于是, 一方面,
    \begin{align*}
        x =
        6\begin{bmatrix}
             1 \\ 0 \\ 0 \\
         \end{bmatrix}
        + 2\begin{bmatrix}
               0 \\ 1 \\ 0 \\
           \end{bmatrix}
        + 8\begin{bmatrix}
               0 \\ 0 \\ 1 \\
           \end{bmatrix}.
    \end{align*}

    另一方面, 我们设
    \begin{align*}
        x =
        k_1 \begin{bmatrix}
                1 \\ 0 \\ 0 \\
            \end{bmatrix}
        + k_2 \begin{bmatrix}
                  0 \\ 1 \\ 0 \\
              \end{bmatrix}
        + k_3 \begin{bmatrix}
                  0 \\ 0 \\ 1 \\
              \end{bmatrix}.
    \end{align*}
    则
    \begin{align*}
        \begin{bmatrix}
            6 \\ 2 \\ 8 \\
        \end{bmatrix}
        =
        \begin{bmatrix}
            k_1 \\ 0 \\ 0 \\
        \end{bmatrix}
        + \begin{bmatrix}
              0 \\ k_2 \\ 0 \\
          \end{bmatrix}
        + \begin{bmatrix}
              0 \\ 0 \\ k_3 \\
          \end{bmatrix}
        = \begin{bmatrix}
              k_1 \\ k_2 \\ k_3 \\
          \end{bmatrix}.
    \end{align*}
    由此可知,
    \(k_1\), \(k_2\), \(k_3\)
    分别是
    \(6\), \(2\), \(8\).
\end{example}

我们可用完全类似的方法证明如下结果.
我留它为您的习题.

\begin{theorem}
    可以, 且只能以一个方式写一个
    \(1 \times n\)~阵为 \(n\)~级单位阵 \(I_n\)~的%
    行的数乘的和.
\end{theorem}

\section{行列式的性质}

本节, 我们学习行列式的几条重要的性质.

\begin{theorem}[规范性]
    单位阵的行列式为 \(1\).
\end{theorem}

\begin{proof}
    我们用数学归纳法证明此事.
    具体地, 设 \(P(n)\) 为命题
    \begin{align*}
        \det {(I_n)} = 1.
    \end{align*}
    则, 我们的目标是,
    对任何正整数 \(n\), \(P(n)\) 是对的.

    \(P(1)\) 显然是对的.
    我就不解释它了; 您解释它.

    现在, 我们假定 \(P(m-1)\) 是对的.
    我们要证 \(P(m)\) 也是对的;
    也就是, \(\det {(I_m)} = 1\).
    按列~\(1\) 展开, 有
    \begin{align*}
             & \det {(I_m)}
        \\
        = {} & \sum_{i = 1}^{m}
        {(-1)^{i+1} [I_m]_{i,1} \det {(I_m (i|1))}}
        \\
        % = {} & \sum_{i = 1}^{m}
        % {(-1)^{i+1} \delta(i, 1) \det {(I_m (i|1))}}
        % \\
        = {} & (-1)^{1+1} \delta(1, 1) \det {(I_m (1|1))}
        + \sum_{i = 2}^{m}
        {(-1)^{i+1} \delta(i, 1) \det {(I_m (i|1))}}
        \\
        = {} & \det {(I_{m-1})}
        + \sum_{i = 2}^{m}
        {(-1)^{i+1} 0 \det {(I_m (i|1))}}
        \\
        = {} & \det {(I_{m-1})}
        \\
        = {} & 1.
    \end{align*}
    所以, \(P(m)\) 是对的.
    由数学归纳法, 待证命题成立.
\end{proof}

我接下来要讲的二个性质的论证利用%
按 (任何) 一列展开行列式%
的公式.
请允许我引用此公式:

% subfile fix: repeat restatable
% see:
% determinanto - determinantoj.tex
\ifSubfilesClassLoaded{%
    \begin{restatable}{theorem}{TheoremExpansionAboutAnyColumn}
        设 \(A\) 是 \(n\)~级阵 (\(n \geq 1\)).
        设 \(j\) 为整数, 且 \(1 \leq j \leq n\).
        则
        \begin{align*}
            \det {(A)} = \sum_{i = 1}^{n}
            {(-1)^{i+j} [A]_{i,j} \det {(A(i|j))}}.
        \end{align*}
    \end{restatable}
}{\TheoremExpansionAboutAnyColumn*}

\begin{theorem}[多线性]
    行列式 (关于列) 是多线性的.
    具体地, 对任何不超过 \(n\) 的正整数 \(j\),
    任何 \(n-1\)~个 \(n \times 1\)~阵
    \(a_1\), \(\dots\), \(a_{j-1}\),
    \(a_{j+1}\), \(\dots\), \(a_n\),
    任何二个 \(n \times 1\)~阵 \(x\), \(y\),
    任何二个数 \(s\), \(t\),
    有
    \begin{align*}
             & \det
        {[a_1, \dots, a_{j-1}, sx + ty, a_{j+1}, \dots, a_n]}
        \\
        = {} &
        s
        \det {[a_1, \dots, a_{j-1}, x, a_{j+1}, \dots, a_n]}
        +
        t
        \det {[a_1, \dots, a_{j-1}, y, a_{j+1}, \dots, a_n]}.
    \end{align*}
    (若 \(j = 1\), 则 \(a_1\) 不出现;
    若 \(j = n\), 则 \(a_n\) 不出现.
    下同.)
\end{theorem}

\begin{proof}
    作三个 \(n\)~级阵 \(A\), \(B\), \(C\):
    \(A\), \(B\), \(C\)~的列~\(k\) 为 \(a_k\) (\(k \neq j\));
    \(A\)~的列~\(j\) 为 \(x\);
    \(B\)~的列~\(j\) 为 \(y\);
    \(C\)~的列~\(j\) 为 \(sx + ty\).
    则
    \begin{align*}
         & \det {(A)}
        = \det
        {[a_1, \dots, a_{j-1}, x, a_{j+1}, \dots, a_n]};
        \\
         & \det {(B)}
        = \det
        {[a_1, \dots, a_{j-1}, y, a_{j+1}, \dots, a_n]};
        \\
         & \det {(C)}
        = \det
        {[a_1, \dots, a_{j-1}, sx + ty, a_{j+1}, \dots, a_n]}.
    \end{align*}
    不难发现, 若 \(k \neq j\), 则
    \(
    [A]_{i,k} = [B]_{i,k} = [C]_{i,k},
    \)
    故
    \(
    A(i|j) = B(i|j) = C(i|j).
    \)
    再注意,
    \(
    [C]_{i,j} = s[A]_{i,j} + t[B]_{i,j},
    \)
    得
    \begin{align*}
             & \det {(C)}
        \\
        = {} &
        \sum_{i = 1}^{n} {
        (-1)^{i+j} [C]_{i,j} \det {(C(i|j))}
        }
        \\
        = {} &
        \sum_{i = 1}^{n} {
        (-1)^{i+j} (s[A]_{i,j} + t[B]_{i,j}) \det {(C(i|j))}
        }
        \\
        = {} &
        s\sum_{i = 1}^{n} {
        (-1)^{i+j} [A]_{i,j} \det {(C(i|j))}
        }
        +
        t\sum_{i = 1}^{n} {
        (-1)^{i+j} [B]_{i,j} \det {(C(i|j))}
        }
        \\
        = {} &
        s\sum_{i = 1}^{n} {
        (-1)^{i+j} [A]_{i,j} \det {(A(i|j))}
        }
        +
        t\sum_{i = 1}^{n} {
        (-1)^{i+j} [B]_{i,j} \det {(B(i|j))}
        }
        \\
        = {} &
        s \det{(A)} + t \det{(B)}.
        \qedhere
    \end{align*}
\end{proof}

多线性有一个直白的推广.
具体地,
对任何不超过 \(n\) 的正整数 \(j\),
任何 \(n-1\)~个 \(n \times 1\)~阵
\(a_1\), \(\dots\), \(a_{j-1}\),
\(a_{j+1}\), \(\dots\), \(a_n\),
任何 \(\ell\)~个 \(n \times 1\)~阵
\(b_1\), \(b_2\), \(\dots\), \(b_\ell\),
任何 \(\ell\)~个数
\(s_1\), \(s_2\), \(\dots\), \(s_\ell\),
有
\begin{align*}
         & \det
    {\left[
            a_1, \dots, a_{j-1},
            \sum_{k = 1}^{\ell} {s_k b_k},
            a_{j+1}, \dots, a_n
            \right]}
    \\
    = {} &
    \sum_{k = 1}^{\ell}
    {s_k
    \det {[a_1, \dots, a_{j-1},
                b_k,
                a_{j+1}, \dots, a_n]}}.
\end{align*}
您可施数学归纳法于 \(\ell\) 以证此事
(注意,
\(
s_1 b_1 + \dots + s_\ell b_\ell
= 1 (s_1 b_1 + \dots + s_{\ell-1} b_{\ell-1}) + s_\ell b_\ell
\)%
);
其实, 这就是多次用多线性的结果.

\begin{theorem}[交错性]
    行列式 (关于列) 是交错性的.
    具体地,
    若 \(n\)~级阵 \(A\) 有二列相同,
    则 \(\det {(A)} = 0\).
\end{theorem}

有一件事值得提.
明显地, \(1\)~级阵不可能有二列相同.
不过,
根据
``若 \(p\) 是错的, 则 `若 \(p\), 则 \(q\)' 是对的''
原则,
我们仍认为 \(1\)~级阵的行列式有交错性.
(``若 \(p\), 则 \(q\)''
相当于
``对适合条件 \(p\) 的任何对象~\(o\),
\(o\) 也适合条件 \(q\)''.
形如 ``对任何 A, 必 B'' 的话%
的反面就是 ``存在某个 A, 使它非 B''.
所以, 若我们找不到 ``非 B'' 的 A,
我们应认为, ``对任何 A, 必 B'' 是对的.
特别地, 当 A 不存在时, 自然没有非 B 的 A.
故, 我们认为, 即使 A 不存在,
``对任何 A, 必 B'' 是对的.)

\begin{proof}
    我们用数学归纳法证明此事.
    具体地, 设 \(P(n)\) 为命题
    \begin{quotation}
        对每一个有相同的二列的 \(n\)~级阵 \(A\),
        它的行列式必为 \(0\).
    \end{quotation}
    则, 我们的目标是,
    对任何正整数 \(n\), \(P(n)\) 是对的.

    前面解释过, 我们认为, \(P(1)\) 是对的.

    现在考虑 \(P(2)\).
    为此, 任取有相同的二列的 \(2\)~级阵
    \begin{align*}
        A = \begin{bmatrix}
                a & a \\
                b & b \\
            \end{bmatrix}.
    \end{align*}
    则
    \begin{align*}
        \det {(A)}
        = ab - ba
        = 0.
    \end{align*}
    故 \(P(2)\) 是对的.

    现在, 我们假定 \(P(m-1)\) 是对的
    (注意, 此处 \(m \geq 3\)).
    我们要证 \(P(m)\) 也是对的.
    任取一个有相同的二列的 \(m\)~级阵 \(A\).
    我们设 \(A\)~的列~\(p\) 等于列~\(q\),
    其中, \(1 \leq p < q \leq m\).
    在 \(1\), \(2\), \(\dots\), \(m\) 这 \(m\)~个数里,
    我们必定能找到一个数 \(j\),
    它既不等于 \(p\), 也不等于 \(q\).
    按列~\(j\) 展开行列式, 有
    \begin{align*}
        \det {(A)}
        = \sum_{i = 1}^{m} {(-1)^{i + j} [A]_{i,j}
        \det {(A(i|j))}}.
    \end{align*}
    注意, \(m-1\)~级阵 \(A(i|j)\) 仍有相同的二列.
    所以, 根据假定, \(\det {(A(i|j))} = 0\).
    故 \(\det {(A)} = 0\).

    所以, \(P(m)\) 是对的.
    由数学归纳法, 待证命题成立.
\end{proof}

行列式的一些性质是多线性与交错性的推论,
如

\begin{theorem}[反称性]
    行列式 (关于列) 是反称性的.
    具体地,
    设 \(A\) 是 \(n\)~级阵.
    设交换 \(A\)~的列~\(p\) 与列~\(q\) 后得到的阵为 \(B\)
    (\(p < q\)).
    则 \(\det {(B)} = -\det {(A)}\).
    (通俗地, 交换方阵的二列, 则行列式变号.)
\end{theorem}

\begin{proof}
    设 \(A = [a_1, a_2, \dots, a_n]\).
    为方便, 我们写
    \begin{align*}
        f(x, y)
        = \det {[a_1, \dots, a_{p-1}, x, a_{p+1}, \dots,
                    a_{q-1}, y, a_{q+1}, \dots, a_n]}.
    \end{align*}
    则
    \begin{align*}
        0
        = {} & f(a_p + a_q, a_p + a_q)               \\
        = {} & f(a_p, a_p + a_q) + f(a_q, a_p + a_q) \\
        = {} & (f(a_p, a_p) + f(a_p, a_q))
        + (f(a_q, a_p) + f(a_q, a_q))                \\
        = {} & f(a_p, a_q) + f(a_q, a_p)             \\
        = {} & \det {(A)} + \det {(B)}.
        \qedhere
    \end{align*}
\end{proof}

\section{用行列式的性质确定行列式}

本节, 我们研究如何%
用行列式的一些性质确定行列式.

假定, 您了解一些人
(而不只是知道他们的存在).
自然地, 您也知道他们的一些特征.
我从这些人里选一个.
我想, 您可以说出此人的几个特征.
不过, 反过来,
若我说出某人的几个特征,
您能确定此人吗?
那一般是不一定的.
不过, 若我给出某些特别的特征,
那是有可能的.

抽象地, \(n\)~级阵的行列式就是一个%
定义在全体 \(n\)~级阵上的函数.
就像一个人有多的特征那样,
行列式也有多的特征 (或者, 性质),
无论是有名的 (有名字的),
还是无名的.
注意,
一些特征并不是行列式特有的:
比如, 零函数 \(z(A) = 0\)
(其中, \(A\) 是任何的 \(n\)~级阵)
适合多线性、交错性、反称性,
且恒取 \(1\) 的函数 \(u(A) = 1\)
适合规范性,
但它们都不是行列式.
不过, 若我们联合行列式的几个特征,
则它们可确定行列式.

\begin{theorem}
    设定义在全体 \(n\)~级阵上的函数 \(f\) 适合:

    (1)
    (多线性)
    对任何不超过 \(n\) 的正整数 \(j\),
    任何 \(n-1\)~个 \(n \times 1\)~阵
    \(a_1\), \(\dots\), \(a_{j-1}\),
    \(a_{j+1}\), \(\dots\), \(a_n\),
    任何二个 \(n \times 1\)~阵 \(x\), \(y\),
    任何二个数 \(s\), \(t\),
    有
    \begin{align*}
             & f
            {([a_1, \dots, a_{j-1}, sx + ty,
                        a_{j+1}, \dots, a_n])}
        \\
        = {} &
        s
        f {([a_1, \dots, a_{j-1}, x, a_{j+1}, \dots, a_n])}
        +
        t
        f {([a_1, \dots, a_{j-1}, y, a_{j+1}, \dots, a_n])}.
    \end{align*}

    (2)
    (交错性)
    若 \(n\)~级阵 \(A\) 有二列相同,
    则 \(f {(A)} = 0\).

    那么, 对任何 \(n\)~级阵 \(A\),
    \(f(A) = f(I) \det {(A)}\).

    特别地, 若 \(f(I) = 1\) (规范性),
    则 \(f\) 就是行列式.
\end{theorem}

\begin{proof}
    我们设 \(g(A) = f(A) - f(I) \det {(A)}\).
    不难验证, \(g\) 也有多线性、交错性、反称性.
    (具体地, 设交换 \(n\)~级阵 \(A\)~的%
    列~\(p\) 与列~\(q\) 后得到的阵为 \(B\),
    且 \(p < q\),
    则 \(g(B) = -g(A)\);
    这是前二个性质的推论.)
    不难看出, \(g(I) = 0\).
    我们的目标是, 证 \(g\) 是零函数 (即, 恒为 \(0\)).

    任取一个 \(n\)~级阵 \(A\).
    设 \(n\)~级单位阵的列~\(1\), \(2\), \(\dots\), \(n\)
    分别是 \(e_1\), \(e_2\), \(\dots\), \(e_n\).
    设 \(A\) 的列~\(1\), \(2\), \(\dots\), \(n\)
    分别是 \(a_1\), \(a_2\), \(\dots\), \(a_n\).
    于是
    \begin{align*}
        a_{k}
        % = {} &
            =
            [A]_{1,k} e_{1} + [A]_{2,k} e_{2}
        + \dots + [A]_{n,k} e_{n}
        % \\
        % = {} &
        =
        \sum_{i_k = 1}^{n} {[A]_{i_k,k} e_{i_k}}.
    \end{align*}
    % (注意, 这里, 我用了 \(n\)~个求和指标
    % \(i_1\), \(i_2\), \(\dots\), \(i_n\).)
    从而, 利用多线性,
    \begin{align*}
        g(A)
        = {} &
        g([a_1, a_2, \dots, a_n])
        \\
        = {} &
        g\left(\left[ \sum_{i_1 = 1}^{n} {[A]_{i_1,1} e_{i_1},
        a_2, \dots, a_n} \right]\right)
        \\
        = {} &
        \sum_{i_1 = 1}^{n} {[A]_{i_1,1}\,
        g([ e_{i_1}, a_2, \dots, a_n ])}
        \\
        = {} &
        \sum_{i_1 = 1}^{n} {[A]_{i_1,1}\,
        g\left(\left[ e_{i_1},
            \sum_{i_2 = 1}^{n} [A]_{i_2,2} e_{i_2},
        \dots, a_n
        \right]\right)}
        \\
        = {} &
        \sum_{i_1 = 1}^{n} {
        \sum_{i_2 = 1}^{n} {[A]_{i_1,1} [A]_{i_2,2}\,
        g([ e_{i_1}, e_{i_2}, \dots, a_n ])}}
        \\
        = {} &
        \dots
        \\
        = {} &
        \sum_{i_1 = 1}^{n} {
        \sum_{i_2 = 1}^{n} {
        \dots
        \sum_{i_n = 1}^{n} {
        [A]_{i_1,1} [A]_{i_2,2} \dots [A]_{i_n,n} \,
        g([e_{i_1}, e_{i_2}, \dots, e_{i_n}])}}}.
    \end{align*}

    利用交错性, 不难看出,
    若存在整数 \(p\), \(q\)
    使 \(1 \leq p < q \leq n\),
    且 \(i_p = i_q\),
    则 \([e_{i_1}, e_{i_2}, \dots, e_{i_n}]\)
    有二列相同,
    故 \(g([e_{i_1}, e_{i_2}, \dots, e_{i_n}]) = 0\).
    所以,
    \begin{align*}
        g(A)
        = \sum_{\substack{
        1 \leq i_1, i_2, \dots, i_n \leq n \\
                i_1, i_2, \dots, i_n\,\text{互不相同}
            }}
        {[A]_{i_1,1} [A]_{i_2,2} \dots [A]_{i_n,n}\,
            g([e_{i_1}, e_{i_2}, \dots, e_{i_n}])}.
    \end{align*}

    我们用反称性证明每一个
    \(g([e_{i_1}, e_{i_2}, \dots, e_{i_n}])\)
    (其中, \(i_1\), \(i_2\), \(\dots\), \(i_n\)
    是不超过 \(n\) 的正整数, 且互不相同)
    都是 \(0\).
    适当地交换
    \([e_{i_1}, e_{i_2}, \dots, e_{i_n}]\)~的列,
    可变它为 \([e_1, e_2, \dots, e_n]\),
    即 \(n\)~级单位阵.
    从而
    \begin{align*}
        g([e_{i_1}, e_{i_2}, \dots, e_{i_n}])
        = \pm g(I) = 0.
    \end{align*}

    所以 \(g(A) = 0\).
    故 \(f(A) = f(I) \det {(A)}\).
    (反过来, 不难验证, 若我们定义
    \(f(A) = f(I) \det {(A)}\),
    则 \(f\) 适合多线性与交错性.)
\end{proof}

\begin{theorem}
    若定义在全体 \(n\)~级阵上的函数 \(f\)
    适合规范性、多线性、交错性,
    则 \(f\) 就是 (\(n\)~级阵的) 行列式 (函数).
\end{theorem}

由此, 理论地, 我们可用%
行列式的这三条性质%
定义行列式:

\begin{definition}
    设 \(f\) 是定义在全体 \(n\)~级阵上的函数.
    若 \(f\) 适合如下三条, 则说
    \(f\) 是 (\(n\)~级阵的) 一个行列式函数
    (自然地, 若 \(A\) 是 \(n\)~级阵,
    则 \(f(A)\) 是 \(A\) 的一个行列式):

    (1)
    (规范性)
    若 \(I\) 是 \(n\)~级单位阵,
    则 \(f(I) = 1\).

    (2)
    (多线性)
    对任何不超过 \(n\) 的正整数 \(j\),
    任何 \(n-1\)~个 \(n \times 1\)~阵
    \(a_1\), \(\dots\), \(a_{j-1}\),
    \(a_{j+1}\), \(\dots\), \(a_n\),
    任何二个 \(n \times 1\)~阵 \(x\), \(y\),
    任何二个数 \(s\), \(t\),
    有
    \begin{align*}
             & f
            {([a_1, \dots, a_{j-1}, sx + ty,
                        a_{j+1}, \dots, a_n])}
        \\
        = {} &
        s
        f {([a_1, \dots, a_{j-1}, x, a_{j+1}, \dots, a_n])}
        +
        t
        f {([a_1, \dots, a_{j-1}, y, a_{j+1}, \dots, a_n])}.
    \end{align*}

    (3)
    (交错性)
    若 \(n\)~级阵 \(A\) 有二列相同,
    则 \(f {(A)} = 0\).
\end{definition}

不过, 我没有这么作,
因为, 若我们如此定义行列式,
则我们要证明 (\(n\)~级阵的) 行列式函数存在且唯一,
而这可能是难的.

\section{\texorpdfstring{``行'' 列式}%
  {"行" 列式}}

前面, 我们学习了行列式的一些公式与性质.
它们至少有一个共同点:
它们都是关于列的命题.
行列式的定义就是按列~\(1\) 展开.
利用定义,
我们得到了按任何一列展开行列式的公式.
然后, 我们得到了行列式的一些关于列的性质.
自然地, 我们问:
行列式是否也有关于行的公式与性质?
此事的回答是 ``是''.
为此, 我们先按行~\(1\) 展开行列式.

% restatable
\begin{restatable}{theorem}{TheoremExpansionAboutRowOne}
    设 \(A\) 是 \(n\)~级阵 (\(n \geq 1\)).
    则
    \begin{align*}
        \det {(A)} = \sum_{j = 1}^{n}
        {(-1)^{1+j} [A]_{1,j} \det {(A(1|j))}}.
    \end{align*}
\end{restatable}

我们无妨先用小级阵验证此命题.

设 \(A\) 是 \(1\)~级阵.
显然
(我作过一个关于 ``\(0\)~级阵'' 及它的行列式的约定).

设 \(A\) 是 \(2\)~级阵 (也就是, 取 \(n = 2\)).
则
\begin{align*}
    \det {(A)}
    = {} &
    [A]_{1,1} [A]_{2,2} - [A]_{2,1} [A]_{1,2}
    \\
    = {} &
    [A]_{1,1} [A]_{2,2} - [A]_{1,2} [A]_{2,1}
    \\
    = {} &
    (-1)^{1+1} [A]_{1,1} \det {(A(1|1))}
    + (-1)^{1+2} [A]_{1,2} \det {(A(1|2))}.
\end{align*}
所以, \(n = 2\) 时, 命题是对的.

\begin{proof}
    我们用数学归纳法证明此事.
    具体地, 设 \(P(n)\) 为命题
    \begin{quotation}
        对任何 \(n\)~级阵 \(A\),
        \begin{align*}
            \det {(A)} = \sum_{j = 1}^{n}
            {(-1)^{1+j} [A]_{1,j} \det {(A(1|j))}}.
        \end{align*}
    \end{quotation}
    则, 我们的目标是,
    对任何正整数 \(n\), \(P(n)\) 是对的.

    我们已知 \(P(1)\), \(P(2)\) 都是对的.

    现在, 我们假定 \(P(m-1)\) 是对的.
    我们要证 \(P(m)\) 也是对的.
    任取一个 \(m\)~级阵 \(A\).
    为方便, 我们写
    \((-1)^{1+1} [A]_{1,1} \det {(A(1|1))}\)
    为 \(f\).
    于是%
    \begingroup
    \reqnomode
    \begin{align*}
             & \det {(A)} \\
        = {} &
        \sum_{i = 1}^{m} {
        (-1)^{i+1} [A]_{i,1} \det {(A(i|1))}
        }
        \tag*{(1)}
        \\
        = {} &
        f
        +
        \sum_{i = 2}^{m} {
        (-1)^{i+1} [A]_{i,1} \det {(A(i|1))}
        }
        \\
        = {} &
        f
        +
        \sum_{i = 2}^{m} {
        (-1)^{i+1} [A]_{i,1}
        \sum_{j = 2}^{m} {
        (-1)^{1+j-1} [A]_{1,j} \det {(A(i,1|1,j))}
        }
        }
        \tag*{(2)}
        \\
        = {} &
        f
        +
        \sum_{i = 2}^{m} {
        \sum_{j = 2}^{m} {
        (-1)^{i+1} [A]_{i,1}
        (-1)^{1+j-1} [A]_{1,j} \det {(A(i,1|1,j))}
        }
        }
        \tag*{(3)}
        \\
        = {} &
        f
        +
        \sum_{j = 2}^{m} {
        \sum_{i = 2}^{m} {
        (-1)^{i+1} [A]_{i,1}
        (-1)^{1+j-1} [A]_{1,j} \det {(A(i,1|1,j))}
        }
        }
        \tag*{(4)}
        \\
        = {} &
        f
        +
        \sum_{j = 2}^{m} {
        \sum_{i = 2}^{m} {
        (-1)^{1+j} [A]_{1,j}
        (-1)^{i-1+1} [A]_{i,1} \det {(A(i,1|1,j))}
        }
        }
        \tag*{(5)}
        \\
        = {} &
        f
        +
        \sum_{j = 2}^{m} {
        (-1)^{1+j} [A]_{1,j}
        \sum_{i = 2}^{m} {
        (-1)^{i-1+1} [A]_{i,1} \det {(A(i,1|1,j))}
        }
        }
        \tag*{(6)}
        \\
        = {} &
        f
        +
        \sum_{j = 2}^{m} {
        (-1)^{1+j} [A]_{1,j} \det {(A(1|j))}
        }
        \tag*{(7)}
        \\
        = {} &
        \sum_{j = 1}^{m} {
        (-1)^{1+j} [A]_{1,j} \det {(A(1|j))}
        }.
    \end{align*}%
    \endgroup
    所以, \(P(m)\) 是对的.
    由数学归纳法, 待证命题成立.

    我想, 您看到了公式右侧的序号.
    这是方便我说话用的.
    我想, 适当地解释这几步,
    是有好处的.

    (1) 是行列式的定义.

    (2) 利用了假定.
    设 \(i > 1\).
    我们假定可按行~\(1\) 展开任何 \(m - 1\)~级阵的行列式.
    \(A(i|1)\) 不就是 \(m - 1\)~级阵吗?
    那么, 我们就按 \(A(i|1)\) 的行~\(1\) 展开.
    \(A(i|1)\) 的行~\(1\) 正好对应
    \(A\) 的行~\(1\).
    最后, 注意,
    \(A\)~的 \((1, j)\)-元恰是
    \(A(i|1)\)~的 \((1, j - 1)\)-元.

    (3) 利用了分配律 (还有加法的结合律与交换律).

    (4) 利用了加法的结合律与交换律.
    (通俗地, 求和号的次序可换.)

    (5) 利用了 \(-1\)~的整数次方的性质 (我在前面说过).
    当然, 我还利用了乘法的结合律与交换律.

    (6) 又用了一次分配律 (还有加法的结合律与交换律).
    不过, 跟 (3) 对比, 这次是反过来用.

    (7) 用到了行列式的定义.
    注意, \(i > 1\) 时,
    \(A\)~的 \((i, 1)\)-元恰是
    \(A(1|j)\)~的 \((i - 1, 1)\)-元.
\end{proof}

由此, 我们不难证明行列式与转置的关系.

% 不难看到, 按行~\(1\) 展开行列式的公式%
% 跟按列~\(1\) 展开行列式的公式是类似的.
% 为方便,
% 我作如下定义.
% % 我作一个暂时的定义.

% \begin{definition}[``行'' 列式]
%     设 \(A\) 是 \(n\)~级阵 (\(n \geq 1\)).
%     定义 \(A\) 的 ``行'' 列式
%     \begin{align*}
%         \operatorname{det}' {(A)}
%         =
%         \begin{dcases}
%             [A]_{1,1},
%              & n = 1;    \\
%             \sum_{j = 1}^{n}
%             {(-1)^{1+j} [A]_{1,j}
%             \operatorname{det}' {(A(1|j))}},
%              & n \geq 2.
%         \end{dcases}
%     \end{align*}
% \end{definition}

% 由上个定理,
% 不难用数学归纳法证明,
% ``行'' 列式就是行列式.
% 作命题 \(P(n)\):
% 对任何 \(n\)~级阵 \(A\),
% \(A\)~的 ``行'' 列式等于 \(A\)~的行列式.
% \(P(1)\) 显然是对的.
% 假定 \(P(m-1)\) 是对的.
% 任取一个 \(m\)~级阵 \(A\).
% 则
% \begin{align*}
%     \operatorname{det}' {(A)}
%     = {} & \sum_{j = 1}^{m}
%     {(-1)^{1+j} [A]_{1,j} \operatorname{det}' {(A(1|j))}}
%     \\
%     = {} & \sum_{j = 1}^{m}
%     {(-1)^{1+j} [A]_{1,j} \det {(A(1|j))}}
%     \\
%     = {} & \det {(A)}.
% \end{align*}

% 若我用
% ``行'' 列式~\(\operatorname{det}'\)
% 作行列式的定义,
% 那么, 我可用完全类似的方法,
% 根据新的定义,
% 证明按一行展开行列式的公式
% (见本章, 节~\sekcio{7}):%
% \begingroup
% \reqnomode
% \begin{align*}
%          & \det {(A)}                  \\
%     = {} &
%     \sum_{j = 1}^{m} {
%     (-1)^{1+j} [A]_{1,j} \det {(A(1|j))}
%     }
%     % \tag*{(1)}
%     \\
%     = {} &
%     \sum_{j = 1}^{m} {
%     (-1)^{1+j} [A]_{1,j}
%     \sum_{\substack{1 \leq \ell \leq m \\ \ell \neq j}} {
%     (-1)^{(i - 1) + (\ell - \rho(\ell, j))}
%         [A]_{i,\ell} \det {(A({1,i}|{j,\ell}))}
%     }
%     }
%     % \tag*{(2)}
%     \\
%     = {} &
%     \sum_{j = 1}^{m} {
%     \sum_{\substack{1 \leq \ell \leq m \\ \ell \neq j}} {
%     (-1)^{(i - 1) + (\ell - \rho(\ell, j))}
%     (-1)^{1+j}
%         [A]_{1,j} [A]_{i,\ell} \det {(A({1,i}|{j,\ell}))}
%     }
%     }
%     % \tag*{(3)}
%     \\
%     = {} &
%     \sum_{\ell = 1}^{m} {
%     \sum_{\substack{1 \leq j \leq m    \\ j \neq \ell}} {
%     (-1)^{(i - 1) + (\ell - \rho(\ell, j))}
%     (-1)^{1+j}
%         [A]_{1,j} [A]_{i,\ell} \det {(A({1,i}|{j,\ell}))}
%     }
%     }
%     % \tag*{(4)}
%     \\
%     = {} &
%     \sum_{\ell = 1}^{m} {
%     \sum_{\substack{1 \leq j \leq m    \\ j \neq \ell}} {
%     (-1)^{i} (-1)^{\ell} (-1)^{\rho(\ell, j)} (-1)^{j}
%         [A]_{1,j} [A]_{i,\ell} \det {(A({1,i}|{j,\ell}))}
%     }
%     }
%     % \tag*{(5)}
%     \\
%     = {} &
%     \sum_{\ell = 1}^{m} {
%     \sum_{\substack{1 \leq j \leq m    \\ j \neq \ell}} {
%     (-1)^{i + \ell} (-1)^{1 + j - \rho(j, \ell)}
%         [A]_{1,j} [A]_{i,\ell} \det {(A({1,i}|{j,\ell}))}
%     }
%     }
%     % \tag*{(6)}
%     \\
%     = {} &
%     \sum_{\ell = 1}^{m} {(-1)^{i + \ell} [A]_{i,\ell}
%     \sum_{\substack{1 \leq j \leq m    \\ j \neq \ell}} {
%     (-1)^{1 + j - \rho(j, \ell)}
%         [A]_{1,j} \det {(A({1,i}|{j,\ell}))}
%     }
%     }
%     % \tag*{(7)}
%     \\
%     = {} &
%     \sum_{\ell = 1}^{m} {(-1)^{i + \ell} [A]_{i,\ell}
%     \det {(A(i|\ell))}
%     },
%     % \tag*{(8)}
% \end{align*}%
% \endgroup
% % 不过, 写几乎一样的论证是无趣的.
% % 所以, 换个思路.
% 其中, 对整数 \(i\), \(j\),
% \begin{align*}
%     \rho(i, j)
%     = \begin{cases}
%           0, & i < j;    \\
%           1, & i \geq j.
%       \end{cases}
% \end{align*}
% 不过, 我们换个思路.

\begin{theorem}
    设 \(A\) 是 \(n\)~级阵 (\(n \geq 1\)).
    则 \(A\)~的转置, \(A^{\mathrm{T}}\),
    的行列式等于 \(A\)~的行列式.
\end{theorem}

\begin{proof}
    我们用数学归纳法证明此事.
    具体地, 设 \(P(n)\) 为命题
    \begin{quotation}
        对任何 \(n\)~级阵 \(A\),
        \begin{align*}
            \det {(A^{\mathrm{T}})} = \det {(A)}.
        \end{align*}
    \end{quotation}
    则, 我们的目标是,
    对任何正整数 \(n\), \(P(n)\) 是对的.

    \(P(1)\) 是对的:
    \(1\)~级阵的转置就是自己.

    现在, 我们假定 \(P(m-1)\) 是对的.
    我们要证 \(P(m)\) 也是对的.
    任取一个 \(m\)~级阵 \(A\).
    注意, \([A]_{i,j} = [A^{\mathrm{T}}]_{j,i}\).
    于是, 可以验证,
    \(A(i|k) = (A^{\mathrm{T}}(k|i))^{\mathrm{T}}\).
    则
    \begin{align*}
        \det {(A^{\mathrm{T}})}
        = {} & \sum_{i = 1}^{m} {
        (-1)^{1+i} [A^{\mathrm{T}}]_{1,i}
        \det {(A^{\mathrm{T}}(1|i))}
        }
        \\
        = {} & \sum_{i = 1}^{m} {
        (-1)^{1+i} [A^{\mathrm{T}}]_{1,i}
        \det {((A^{\mathrm{T}}(1|i))^{\mathrm{T}})}
        }
        \\
        = {} & \sum_{i = 1}^{m} {
        (-1)^{i+1} [A]_{i,1}
        \det {(A(i|1))}
        }
        \\
        = {} & \det {(A)}.
    \end{align*}
    所以, \(P(m)\) 是对的.
    由数学归纳法, 待证命题成立.
\end{proof}

此性质是重要的:
利用它,
我们可译行列式的关于列的命题%
为关于行的命题.
我用一个例助您理解;
它是按行~\(i\) 展开行列式的公式.

% restatable that one has manually to repeat
% in at least one other file:
% determinanto - matrica multipliko.tex
\begin{restatable}{theorem}{TheoremExpansionAboutAnyRow}
    设 \(A\) 是 \(n\)~级阵 (\(n \geq 1\)).
    设 \(i\) 为整数, 且 \(1 \leq i \leq n\).
    则
    \begin{align*}
        \det {(A)} = \sum_{j = 1}^{n}
        {(-1)^{i+j} [A]_{i,j} \det {(A(i|j))}}.
    \end{align*}
\end{restatable}

\begin{proof}
    基本的思想:
    先转置,
    施关于列的命题于转置,
    再转置回来.
    回想,
    \(A^{\mathrm{T}}\)~的列~\(i\)
    跟 \(A\)~的行~\(i\) 对应,
    故我们按列~\(i\) 展开 \(A^{\mathrm{T}}\)~的行列式.
    则
    \begin{align*}
        \det {(A)}
        = {} & \det {(A^{\mathrm{T}})}
        \\
        = {} & \sum_{j = 1}^{n}
        {(-1)^{j+i} [A^{\mathrm{T}}]_{j,i}
        \det {(A^{\mathrm{T}}(j|i))}
        }
        \\
        = {} & \sum_{j = 1}^{n}
        {(-1)^{j+i} [A]_{i,j}
        \det {((A(i|j))^{\mathrm{T}})}
        }
        \\
        = {} & \sum_{j = 1}^{n}
        {(-1)^{i+j} [A]_{i,j}
        \det {(A(i|j))}
        }.
        \qedhere
    \end{align*}
\end{proof}

为方便,
我译前面的行列式的关于列的性质%
为关于行的性质.
不过, 我不证它们.
我留它们为您的习题.

\begin{theorem}[多线性]
    行列式 (关于行) 是多线性的.
    具体地, 对任何不超过 \(n\) 的正整数 \(i\),
    任何 \(n-1\)~个 \(1 \times n\)~阵
    \(a_1\), \(\dots\), \(a_{i-1}\),
    \(a_{i+1}\), \(\dots\), \(a_n\),
    任何二个 \(1 \times n\)~阵 \(x\), \(y\),
    任何二个数 \(s\), \(t\),
    有
    \begin{align*}
        \det
        {\begin{bmatrix}
                 a_1     \\
                 \vdots  \\
                 a_{i-1} \\
                 sx + ty \\
                 a_{i+1} \\
                 \vdots  \\
                 a_{n}   \\
             \end{bmatrix}}
        =
        s
        \det
        {\begin{bmatrix}
                 a_1     \\
                 \vdots  \\
                 a_{i-1} \\
                 x       \\
                 a_{i+1} \\
                 \vdots  \\
                 a_{n}   \\
             \end{bmatrix}}
        +
        t
        \det
        {\begin{bmatrix}
                 a_1     \\
                 \vdots  \\
                 a_{i-1} \\
                 y       \\
                 a_{i+1} \\
                 \vdots  \\
                 a_{n}   \\
             \end{bmatrix}}.
    \end{align*}
    (若 \(i = 1\), 则 \(a_1\) 不出现;
    若 \(i = n\), 则 \(a_n\) 不出现.
    下同.)
\end{theorem}

\begin{theorem}[交错性]
    行列式 (关于行) 是交错性的.
    具体地,
    若 \(n\)~级阵 \(A\) 有二行相同,
    则 \(\det {(A)} = 0\).
\end{theorem}

\begin{theorem}[反称性]
    行列式 (关于行) 是反称性的.
    具体地,
    设 \(A\) 是 \(n\)~级阵.
    设交换 \(A\)~的行~\(p\) 与行~\(q\) 后得到的阵为 \(B\)
    (\(p < q\)).
    则 \(\det {(B)} = -\det {(A)}\).
    (通俗地, 交换方阵的二行, 则行列式变号.)
\end{theorem}

\begin{theorem}
    设定义在全体 \(n\)~级阵上的函数 \(f\) 适合:

    (1)
    (多线性)
    对任何不超过 \(n\) 的正整数 \(i\),
    任何 \(n-1\)~个 \(1 \times n\)~阵
    \(a_1\), \(\dots\), \(a_{i-1}\),
    \(a_{i+1}\), \(\dots\), \(a_n\),
    任何二个 \(1 \times n\)~阵 \(x\), \(y\),
    任何二个数 \(s\), \(t\),
    有
    \begin{align*}
        f
        {\left(\begin{bmatrix}
                       a_1     \\
                       \vdots  \\
                       a_{i-1} \\
                       sx + ty \\
                       a_{i+1} \\
                       \vdots  \\
                       a_{n}   \\
                   \end{bmatrix}\right)}
        =
        s
        f
            {\left(\begin{bmatrix}
                               a_1     \\
                               \vdots  \\
                               a_{i-1} \\
                               x       \\
                               a_{i+1} \\
                               \vdots  \\
                               a_{n}   \\
                           \end{bmatrix}\right)}
        +
        t
        f
            {\left(\begin{bmatrix}
                               a_1     \\
                               \vdots  \\
                               a_{i-1} \\
                               y       \\
                               a_{i+1} \\
                               \vdots  \\
                               a_{n}   \\
                           \end{bmatrix}\right)}.
    \end{align*}

    (2)
    (交错性)
    若 \(n\)~级阵 \(A\) 有二行相同,
    则 \(f {(A)} = 0\).

    那么, 对任何 \(n\)~级阵 \(A\),
    \(f(A) = f(I) \det {(A)}\).

    特别地, 若 \(f(I) = 1\) (规范性),
    则 \(f\) 就是行列式.
\end{theorem}

\begin{theorem}
    若定义在全体 \(n\)~级阵上的函数 \(f\)
    适合规范性、(关于行的) 多线性、(关于行的) 交错性,
    则 \(f\) 就是 (\(n\)~级阵的) 行列式 (函数).
\end{theorem}

由此, 理论地, 我们也可如此%
定义行列式:

\begin{definition}
    设 \(f\) 是定义在全体 \(n\)~级阵上的函数.
    若 \(f\) 适合如下三条, 则说
    \(f\) 是 (\(n\)~级阵的) 一个行列式函数
    (自然地, 若 \(A\) 是 \(n\)~级阵,
    则 \(f(A)\) 是 \(A\) 的一个行列式):

    (1)
    (规范性)
    若 \(I\) 是 \(n\)~级单位阵,
    则 \(f(I) = 1\).

    (2)
    (多线性)
    对任何不超过 \(n\) 的正整数 \(i\),
    任何 \(n-1\)~个 \(1 \times n\)~阵
    \(a_1\), \(\dots\), \(a_{i-1}\),
    \(a_{i+1}\), \(\dots\), \(a_n\),
    任何二个 \(1 \times n\)~阵 \(x\), \(y\),
    任何二个数 \(s\), \(t\),
    有
    \begin{align*}
        f
        {\left(\begin{bmatrix}
                       a_1     \\
                       \vdots  \\
                       a_{i-1} \\
                       sx + ty \\
                       a_{i+1} \\
                       \vdots  \\
                       a_{n}   \\
                   \end{bmatrix}\right)}
        =
        s
        f
            {\left(\begin{bmatrix}
                               a_1     \\
                               \vdots  \\
                               a_{i-1} \\
                               x       \\
                               a_{i+1} \\
                               \vdots  \\
                               a_{n}   \\
                           \end{bmatrix}\right)}
        +
        t
        f
            {\left(\begin{bmatrix}
                               a_1     \\
                               \vdots  \\
                               a_{i-1} \\
                               y       \\
                               a_{i+1} \\
                               \vdots  \\
                               a_{n}   \\
                           \end{bmatrix}\right)}.
    \end{align*}

    (3)
    (交错性)
    若 \(n\)~级阵 \(A\) 有二行相同,
    则 \(f {(A)} = 0\).
\end{definition}

\KunAsteriskoEnEnhavtabelo
\section{\texorpdfstring{``行'' 列式 (续)}%
  {"行" 列式 (续)}}
\SenAsteriskoEnEnhavtabelo

\maldevigalegajxo

本节, 为方便,
我译前面的行列式的关于列的公式%
为关于行的公式.
我先引用我们见过的二个公式.

\TheoremExpansionAboutRowOne*

\TheoremExpansionAboutAnyRow*

现在, 我给出新的公式.
您用转置译已有的公式即可.
我就不证它们了.
% 当然, 您可以用老方法, 几乎完全一样地论证它们.

\begin{theorem}
    设 \(A\) 是 \(n\)~级阵 (\(n \geq 1\)).
    设 \(k\) 是不超过 \(n\) 的正整数.
    设 \(i_1\), \(i_2\), \(\dots\), \(i_k\) 是%
    不超过 \(n\) 的正整数,
    且 \(i_1 < i_2 < \dots < i_k\).
    则
    \begin{align*}
         &
        \det {(A)}
        = \sum_{1 \leq j_1 < j_2 < \dots < j_k \leq n}
        {\det {\left(
                A\binom{i_1, i_2, \dots, i_k}
                {j_1, j_2, \dots, j_k}
                \right)}}
        \\
         &
        \qquad \qquad \qquad
        \cdot (-1)^{i_1 + i_2 + \dots + i_k
            + j_1 + j_2 + \dots + j_k}
        \det {(A({i_1,i_2,\dots,i_k}|{j_1,j_2,\dots,j_k}))}.
    \end{align*}
\end{theorem}

\begin{theorem}
    设 \(A\) 是 \(n\)~级阵 (\(n \geq 1\)).
    设 \(i_1\), \(i_2\), \(\dots\), \(i_n\) 是%
    不超过 \(n\) 的正整数,
    且互不相同.
    则
    \begin{align*}
             & \det {(A)}                  \\
        = {} &
        \sum_{\substack{
        1 \leq j_1, j_2, \dots, j_n \leq n \\
                j_1, j_2, \dots, j_n\,\text{互不相同}
            }}
        {s(i_1, i_2, \dots, i_n)\,
            s(j_1, j_2, \dots, j_n)\,
            [A]_{i_1,j_1} [A]_{i_2,j_2} \dots [A]_{i_n,j_n}}.
    \end{align*}
    特别地, 取 \(i_1\), \(i_2\), \(\dots\), \(i_n\)
    为 \(1\), \(2\), \(\dots\), \(n\),
    并注意, \(s(1, 2, \dots, n) = 1\),
    得
    \begin{align*}
        \det {(A)}
        = {} &
        \sum_{\substack{
        1 \leq j_1, j_2, \dots, j_n \leq n \\
                j_1, j_2, \dots, j_n\,\text{互不相同}
            }}
        { s(j_1, j_2, \dots, j_n)\,
            [A]_{1,j_1} [A]_{2,j_2} \dots [A]_{n,j_n}}.
    \end{align*}
\end{theorem}

\begin{theorem}
    设 \(n\)~级阵 \(A\) 的行~\(1\), \(2\), \(\dots\), \(n\)
    分别为 \(a_1\), \(a_2\), \(\dots\), \(a_n\).
    设 \(\ell_1\), \(\ell_2\), \(\dots\), \(\ell_n\) 是%
    不超过 \(n\) 的正整数,
    且互不相同.
    则
    \begin{align*}
        \det {\begin{bmatrix}
                      a_{\ell_1} \\
                      a_{\ell_2} \\
                      \vdots     \\
                      a_{\ell_n} \\
                  \end{bmatrix}
        }
        = s(\ell_1, \ell_2, \dots, \ell_n) \det {(A)}.
    \end{align*}
\end{theorem}

\end{document}

This material is about properties of determinants.

Four important properties (or rather, three)
uniquely determine the determinant function,
and three important properties (or rather, two)
uniquely determine constant multiples of
the determinant function.

I define determinants recursively,
using expansion about column 1.
I have proved many results of determinants about columns.
Now I want to translate these results into those about rows.
First, I prove that expansion about row 1 is correct.
Then, I prove that
the determinant of the transpose of a square matrix A
is equal to det (A).
This fact allows one to translate
results about columns
into results about rows.

More information is available in appendix C.
